\subsection{Memorability}

To design an effective password composition policy that balances \textbf{complexity} (entropy) and \textbf{memorability} (ease of understanding), we can formulate a multi-objective optimization framework. This framework integrates measures of both complexity and memorability to derive an optimal policy that maximizes security while ensuring usability for users.

### 1. **Definitions and Notations**

- **Password Structure Set**: Let \( \mathcal{S} = \{ S_1, S_2, \dots, S_m \} \) denote the set of all possible password structures, where each structure \( S_j \) is a sequence of elements from the alphabet \( \Sigma \). For example, \( S_j = (w, w, d) \) represents a structure with two words followed by a digit.

- **Probability of a Structure**: Let \( P(S_j) \) denote the probability of selecting structure \( S_j \) in the password dataset. The structure probabilities satisfy:
  \[
  \sum_{j=1}^{m} P(S_j) = 1
  \]

- **Elements within Structures**: For a given structure \( S_j = (e_1, e_2, \dots, e_k) \), each element \( e_i \in \Sigma \) has an associated probability \( P(e_i \mid S_j) \).

- **Entropy (Complexity)**: The entropy \( H \) of the password dataset measures its complexity and is defined as:
  \[
  H = H(\mathcal{S}) + \sum_{j=1}^{m} P(S_j) H(\mathcal{E} \mid S_j)
  \]
  where:
  \[
  H(\mathcal{S}) = - \sum_{j=1}^{m} P(S_j) \log P(S_j)
  \]
  and
  \[
  H(\mathcal{E} \mid S_j) = - \sum_{e_i \in \Sigma_j} P(e_i \mid S_j) \log P(e_i \mid S_j)
  \]

- **Memorability Measure**: Let \( M \) denote the memorability of the password dataset. We define \( M \) as a function that quantifies the ease with which users can remember their passwords. For mathematical modeling, \( M \) can be represented as:
  \[
  M = \sum_{j=1}^{m} P(S_j) M(S_j)
  \]
  where \( M(S_j) \) is the memorability score of structure \( S_j \). Higher values of \( M \) indicate greater memorability.

### 2. **Multi-Objective Optimization Framework**

To balance complexity and memorability, we define an objective function that incorporates both \( H \) and \( M \). This can be achieved through a weighted sum approach or by setting constraints.

#### a. **Weighted Sum Approach**

Define a weight parameter \( \lambda \in [0, 1] \) that balances the importance between complexity and memorability. The objective function \( \mathcal{O} \) to be maximized is:
\[
\mathcal{O} = H + \lambda M
\]
Substituting the definitions of \( H \) and \( M \):
\[
\mathcal{O} = \left[ H(\mathcal{S}) + \sum_{j=1}^{m} P(S_j) H(\mathcal{E} \mid S_j) \right] + \lambda \left[ \sum_{j=1}^{m} P(S_j) M(S_j) \right]
\]
The optimization problem becomes:
\[
\max_{\{ P(S_j), P(e_i \mid S_j) \}} \mathcal{O} = H + \lambda M
\]
**Subject to:**
\[
\sum_{j=1}^{m} P(S_j) = 1
\]
\[
\sum_{e_i \in S_j} P(e_i \mid S_j) = 1 \quad \forall j = 1, 2, \dots, m
\]
\[
P(S_j) \geq 0 \quad \forall j
\]
\[
P(e_i \mid S_j) \geq 0 \quad \forall e_i, S_j
\]

#### b. **Constrained Optimization Approach**

Alternatively, we can maximize complexity \( H \) while ensuring that memorability \( M \) meets a minimum threshold \( M_{\text{min}} \):
\[
\max_{\{ P(S_j), P(e_i \mid S_j) \}} H
\]
**Subject to:**
\[
M \geq M_{\text{min}}
\]
\[
\sum_{j=1}^{m} P(S_j) = 1
\]
\[
\sum_{e_i \in S_j} P(e_i \mid S_j) = 1 \quad \forall j
\]
\[
P(S_j) \geq 0 \quad \forall j
\]
\[
P(e_i \mid S_j) \geq 0 \quad \forall e_i, S_j
\]

### 3. **Modeling the Memorability Function \( M(S_j) \)**

The memorability of a structure \( S_j \) can be modeled based on several factors, such as:

- **Simplicity of Structure**: Structures with fewer and more familiar elements (e.g., common words) are generally more memorable.
- **Semantic Relationships**: Structures where elements have meaningful or coherent relationships may enhance memorability.
- **Pattern Regularity**: Structures with regular patterns or repetitions may be easier to remember.

Mathematically, \( M(S_j) \) can be expressed as:
\[
M(S_j) = \alpha \cdot f_{\text{simplicity}}(S_j) + \beta \cdot f_{\text{semantic}}(S_j) + \gamma \cdot f_{\text{pattern}}(S_j)
\]
where \( \alpha, \beta, \gamma \) are weighting coefficients that balance the contribution of each factor, and:
\[
f_{\text{simplicity}}(S_j) = \text{Function quantifying simplicity}
\]
\[
f_{\text{semantic}}(S_j) = \text{Function quantifying semantic coherence}
\]
\[
f_{\text{pattern}}(S_j) = \text{Function quantifying pattern regularity}
\]

### 4. **Incorporating Constraints for Policy Design**

To ensure that the password composition policy is practical and enforceable, we may incorporate additional constraints such as:

- **Minimum and Maximum Structure Length**:
  \[
  k_{\text{min}} \leq |S_j| \leq k_{\text{max}} \quad \forall S_j \in \mathcal{S}
  \]
  where \( |S_j| \) denotes the number of elements in structure \( S_j \).

- **Prohibited Patterns or Elements**:
  \[
  P(e_i \mid S_j) = 0 \quad \text{for prohibited } e_i \text{ in structure } S_j
  \]

- **Minimum Entropy Requirements**:
  \[
  H \geq H_{\text{min}}
  \]

### 5. **Strategy for Designing the Password Composition Policy**

Based on the above framework, the strategy to design the password composition policy involves the following steps:

#### a. **Define the Structure Set \( \mathcal{S} \)**
Determine the allowable password structures based on desired elements (e.g., words, numbers, symbols) and their sequence.

#### b. **Model Complexity and Memorability**
- **Complexity**: Calculate the entropy \( H \) for each structure and the overall dataset.
- **Memorability**: Define and compute the memorability function \( M(S_j) \) for each structure.

#### c. **Formulate the Optimization Problem**
Decide on the objective (e.g., maximize \( H + \lambda M \) or maximize \( H \) subject to \( M \geq M_{\text{min}} \)) and set up the corresponding optimization problem.

#### d. **Solve the Optimization Problem**
Determine the optimal structure probabilities \( P(S_j) \) and element probabilities \( P(e_i \mid S_j) \) that achieve the desired balance between complexity and memorability.

#### e. **Implement Policy Constraints**
Incorporate practical constraints to ensure the policy is user-friendly and enforceable.

#### f. **Iterative Refinement**
Continuously monitor password usage and feedback to adjust the policy parameters \( \lambda \), \( \alpha \), \( \beta \), \( \gamma \), and the structure set \( \mathcal{S} \) to optimize both security and usability.

### 6. **Illustrative Example**

Consider a simplified scenario with two possible structures:
- \( S_1 = (w, w, w) \) (three words)
- \( S_2 = (w, d, s) \) (word, digit, symbol)

Assume the following for memorability:
\[
M(S_1) = 0.8 \quad \text{and} \quad M(S_2) = 0.6
\]

And entropy calculations:
\[
H(\mathcal{S}) = - P(S_1) \log P(S_1) - P(S_2) \log P(S_2)
\]
\[
H(\mathcal{E} \mid S_1) = H_{\text{words}} \quad \text{and} \quad H(\mathcal{E} \mid S_2) = H_{\text{word}} + H_{\text{digit}} + H_{\text{symbol}}
\]

Using the weighted sum approach with \( \lambda = 0.5 \):
\[
\mathcal{O} = H(\mathcal{S}) + P(S_1) H_{\text{words}} + P(S_2) (H_{\text{word}} + H_{\text{digit}} + H_{\text{symbol}}) + 0.5 [P(S_1) \cdot 0.8 + P(S_2) \cdot 0.6]
\]

Solving the optimization problem would yield the optimal probabilities \( P(S_1) \) and \( P(S_2) \) that maximize \( \mathcal{O} \), thereby balancing complexity and memorability.

### 7. **Conclusion and Policy Recommendation**

To design a robust password composition policy for your website that maximizes security while ensuring ease of use:

1. **Integrate Complexity and Memorability**: Utilize a multi-objective optimization framework that considers both entropy (complexity) and a defined memorability measure.

2. **Diversify Structure Probabilities**: Allow a variety of password structures with probabilities optimized to balance entropy and memorability. Avoid enforcing a single structure (e.g., three words) to prevent predictability and potential entropy reduction.

3. **Optimize Structure Selection**: Allocate higher probabilities to structures that offer higher entropy and sufficient memorability, ensuring that users are encouraged to adopt strong yet memorable password formats.

4. **Set Practical Constraints**: Incorporate constraints to maintain usability, such as limiting structure complexity to manageable levels and ensuring that all allowed structures meet minimum security requirements.

5. **Continuous Evaluation and Adjustment**: Regularly assess the effectiveness of the password policy by analyzing password entropy and user feedback on memorability, adjusting the policy parameters as necessary to maintain an optimal balance.

By adopting this strategic, mathematically grounded approach, your website can implement a password composition policy that enhances security through high-entropy passwords while maintaining user-friendly memorability.